\section{Discussion}
\label{sec:discussion}

The results reveal a hierarchy of structures linking the geometry of the
four-primary central configuration to the dynamics of the infinitesimal
body. The central-configuration set determines which primary geometries are
continuously accessible; along those components the restricted equilibria
undergo folds and changes of stability; the resulting centre modes generate
periodic-orbit families; and the associated critical Jacobi levels organize
the energetically accessible regions. We discuss these implications in
turn, together with the limits within which they have been established.


\subsection{Central-configuration topology and equilibrium bifurcation structure}
\label{subsec:discussion_cc}

The principal structural result is that equilibrium multiplicity cannot be
interpreted independently of the component structure of the underlying
central-configuration set. On the fixed slice
\[
(a,b)=(0.9,-0.1),
\]
the two reference configurations with five and nine restricted equilibria
belong to numerically distinct components of the regular zero set
$\psi=0$. Restricting these components to positive masses produces two
different admissible arcs. Consequently, the difference between the
five- and nine-equilibrium reference cases cannot be interpreted as the
result of a single bifurcation sequence connecting the two sampled
configurations.

This observation changes the role of the primary central configuration in
the restricted problem. Rather than serving only as a fixed background
geometry, its component structure determines which equilibrium branches
can coexist along a continuous deformation and which sequences of
bifurcations are accessible on that deformation. On the first admissible
arc the interior multiplicity sequence is
\[
5\rightarrow7\rightarrow5,
\]
before the positive-mass problem approaches a mass-vanishing boundary. On
the second arc the sequence is
\[
7\rightarrow9\rightarrow11\rightarrow9\rightarrow7\rightarrow5.
\]
Thus the equilibrium count is not a monotone measure of deformation or
asymmetry; it depends on the ordering of individual folds along the
admissible component.

The eleven-equilibrium interval illustrates why the scalar count alone is
an incomplete description. The pair $B_a$ created at the
$9\rightarrow11$ fold $F_4$ is not the pair destroyed at the subsequent
$11\rightarrow9$ fold $F_5$. Instead, the pre-existing pair $B_b$ is
destroyed at $F_5$, while $B_a$ persists until $F_6$. The two
nine-equilibrium intervals on either side of the eleven-equilibrium window
therefore have identical multiplicity but different branch populations.
A return to a previous value of $N_{\rm eq}$ does not imply a return to the
same equilibrium configuration.

The mechanisms responsible for changes of multiplicity must also be kept
distinct. The seven interior transitions are bifurcations of the restricted
equilibrium equations with all four primary masses positive. They satisfy
the numerical degeneracy and nondegeneracy conditions for generic
saddle-node folds and exhibit the expected square-root separation law.
By contrast, the Case-1 $5\rightarrow4$ transition occurs as
$\rho_2\rightarrow0$ near the boundary of the admissible primary set. It
is therefore a parameter-boundary degeneration of the four-primary problem
rather than an interior saddle-node bifurcation of the restricted dynamics.

These results suggest a broader organizing principle: in restricted
problems with sufficiently general primary geometries, the component
structure of the admissible primary-configuration set may organize the
accessible restricted bifurcation sequences. The present calculation
supports this interpretation on one two-dimensional slice only; whether it
persists in the full four-parameter admissible set or in other
non-symmetric restricted problems remains an open question.


\subsection{Stability and nonlinear periodic dynamics}
\label{subsec:discussion_dynamics}

The equilibrium continuation also shows that spectral stability is not
restricted to isolated symmetric configurations. Both admissible components
contain finite intervals on which selected libration points are
doubly elliptic. On the reliably tracked Case-1 branch this occurs over
approximately
\[
s\in[0.643,0.835],
\]
while the Case-2 component contains two distinct stable branches over
overlapping parameter ranges. In particular, two different spectrally
stable equilibria coexist near
\[
s\simeq0.196,
\]
one close to a $5{:}2$ centre-frequency commensurability and the other
close to $2{:}1$. This coexistence reinforces the conclusion that neither
the arc coordinate $s$ nor the equilibrium count by itself uniquely
characterizes the local dynamics; equilibrium branch identity is also
required.

The centre modes of the selected equilibria organize the local nonlinear
periodic dynamics. At each parent used in Section~\ref{sec:periodic}, the
selected imaginary eigenvalue satisfies the exact non-resonance condition
required by the Lyapunov centre theorem. The ten computed families are
therefore Lyapunov families in the local theorem sense, including those
continued from saddle$\times$centre parents. What differs is their
stability.

For the selected doubly-elliptic parents, the computed families retain
their nontrivial Floquet multipliers on the unit circle throughout the
continued intervals and are therefore Floquet stable in the sense defined
in Section~\ref{subsec:floquet}. This is a periodic-orbit statement and
should not be interpreted as a proof of nonlinear or KAM stability.
Conversely, families~5--6 originate from saddle$\times$centre parents.
Their centre directions still generate valid Lyapunov families, but the
continued periodic orbits possess a strongly real reciprocal Floquet pair
and are correspondingly real unstable. The existence of a local periodic
family and its Floquet stability are therefore logically distinct
questions.

The low-order centre-frequency near-commensurabilities provide useful
locations at which to examine possible changes in periodic-orbit stability,
but they must not be confused with exact resonances. Families~7--8 are
continued from the Case-1 parent near $3{:}2$, while families~9--10 use the
distinct Case-2 parent near $5{:}2$. In both cases the parent frequency
ratio differs from the exact rational value, so the Lyapunov-centre
non-resonance condition remains satisfied.

The most pronounced nonlinear stability feature occurs on family~10. Its
nontrivial Floquet pair approaches the period-doubling collision at
\[
\nu=-2
\]
to within approximately
\[
|\nu+2|\simeq10^{-7}.
\]
An ultra-fine amplitude sweep with step $8\times10^{-4}$ resolves the
minimum and shows the index turning back without crossing the threshold.
The computed family therefore provides no evidence of a period-doubling
bifurcation over the continuation interval investigated. Family~8 also
approaches the same boundary, but only to within approximately
$2\times10^{-3}$.

This negative result should be interpreted narrowly. The multiplier scans
show no crossing of the selected low-order conditions
\[
\nu=2,\qquad
\nu=0,\qquad
\nu=-1,\qquad
\nu=-2
\]
along the computed portions of the ten families. We have not performed an
independent exhaustive search for all secondary periodic families, nor
does the absence of a crossing on these continuation intervals exclude
period-doubled or resonant solutions elsewhere in parameter or phase space.

The Jacobi and Hill-region calculations provide the geometric context for
these periodic motions. The equilibrium branches determine a ladder of
critical Jacobi values, and changes across these levels can reorganize the
accessible region. The representative Case-2 calculation demonstrates a
net change from one to two connected Hill components across a cluster of
three nearby critical levels. The periodic examples then show selected
Floquet-stable families embedded within the allowed regions at their own
Jacobi constants. The Hill geometry determines where motion at a given
energy is possible, but it does not determine the Floquet stability of the
orbit contained within that region.


\subsection{Relation to previous models, scope and limitations}
\label{subsec:discussion_scope}

The present problem differs structurally from the symmetric restricted
five-body models most commonly treated in the literature. In those models,
reflection, permutation, or rotational symmetry constrains the primary
geometry to a prescribed low-dimensional family. For example, Ollongren's
symmetric restricted five-body configuration supports nine libration
points, including stable equilibria \cite{ollongren1988}, while the ring
problem and other symmetric formulations exploit rotational or reflection
structure to organize their equilibrium and periodic-orbit families
\cite{kalvouridis1999,zotospapadakis2019,zotossuraj2018}. Symmetric
four-body central configurations such as the Caledonian and kite families
similarly restrict the primary configuration to a special geometric locus
\cite{stevesroy1998,benhammouda2020}.

The distinction in the present problem is not merely that the equilibrium
count is larger. On the investigated non-symmetric slice, the admissible
central-configuration set itself separates into different components, and
the restricted equilibrium multiplicity must therefore be interpreted
component by component. The eleven-equilibrium state exceeds the
nine-equilibrium counts reported in the symmetric restricted five-body
examples cited above, but its more significant feature is that it arises
within a continuously resolved saddle-node cascade and contains a branch
exchange. Likewise, the periodic families are continued from
asymmetrically located equilibria using a general shooting formulation and
full-monodromy Floquet analysis rather than being constructed through an
assumed reflection symmetry. The comparison with previous work is therefore
primarily structural rather than a quantitative comparison at matched
masses.

Several limitations delimit these conclusions. First, all primary
continuation is carried out on the fixed slice
\[
(a,b)=(0.9,-0.1).
\]
The numerically distinct components identified there need not remain
disconnected when $a$ and $b$ are allowed to vary. Determining the topology
of the complete admissible set would require continuation in the full
four-dimensional geometric parameter space.

Second, equilibrium completeness is numerical rather than analytic. The
reported multiplicities are supported by deterministic grid searches,
fixed-seed random starts, local searches around the primaries, strict
residual tests, and an expanded re-search of the eleven-equilibrium
configuration. These procedures provide strong numerical evidence for the
reported counts under the stated search protocol, but they do not provide
an analytic upper bound on the number of real equilibria.

Third, branch tracking becomes unreliable near the mass-vanishing limits
where some equilibria recede rapidly. Such points are explicitly flagged
and excluded from branch-dependent stability, centre-frequency, Jacobi,
and periodic-orbit claims. The reliably tracked intervals are therefore
shorter than the full ranges over which individual numerical roots may
still be found.

Fourth, the periodic-orbit calculations use natural amplitude continuation
rather than pseudo-arclength continuation. The reported 198 solutions
therefore represent the portions of ten families accessible while the
chosen amplitude remains a useful continuation coordinate; they are not
claimed to be maximal families. Likewise, no systematic continuation of
secondary resonant or period-doubled branches has been undertaken.

Fifth, the Hill-region calculation is representative rather than
exhaustive. The connectivity change reported in Section~\ref{sec:hill} is
resolution-checked on $300^2$, $400^2$, and $500^2$ grids, but the selected
Jacobi interval spans three nearby equilibrium critical values. The
calculation therefore establishes a net $1\rightarrow2$ connectivity
change across that part of the critical-level ladder rather than assigning
the change to one particular critical equilibrium. A complete topological
classification would require resolving the Hill geometry on both sides of
every relevant critical level.

Finally, the stability results have deliberately limited meaning. We
establish spectral stability of equilibria and Floquet stability of
periodic orbits. We do not establish nonlinear, orbital, or KAM stability,
and finite-time numerical boundedness is not used as a substitute for any
of these notions.

Within these limits, the principal structural conclusions are robust: on
the investigated slice the admissible central-configuration component
structure organizes the accessible equilibrium bifurcation sequences;
those sequences contain generic folds, stable-equilibrium intervals, and a
branch-exchange eleven-equilibrium regime; and selected centre modes extend
to well-resolved Lyapunov periodic families whose Floquet behaviour and
Hill-region setting can be followed quantitatively.